\documentclass[11pt]{tibop-article}

% ============================================================
% PLACEHOLDERS — search for \FILL to find items to complete
% on Overleaf before submission:
%   \FILL{ORCID_AUTHOR1}  — ORCID of first author
%   \FILL{ORCID_AUTHOR2}  — ORCID of second author
%   \FILL{ORCID_AUTHOR3}  — ORCID of third author
%   \FILL{AFF1}           — Affiliation 1 (name, city, country)
%   \FILL{AFF2}           — Affiliation 2 (if different)
%   \FILL{CORR_EMAIL}     — Corresponding author email
%   \FILL{TASK3_METRIC}   — Task 3 (summarization) metric name
%   \FILL{TASK3_SCORE}    — Task 3 final CodaBench score
%   \FILL{TASK4_METRIC}   — Task 4 (VQA) metric name
%   \FILL{TASK4_SCORE}    — Task 4 final CodaBench score
%   \FILL{REPO_DOI}       — DOI of GitHub/Zenodo code repo
%   \FILL{BUNDLE_NAME}    — Confirm exact proceedings bundle name with organizers
% ============================================================
\newcommand{\FILL}[1]{\textbf{[#1]}}

% ---------- Authors ----------
\TIBauthor[VLMinators]{%
  First Author\affOne\orcidlink{\FILL{ORCID\_AUTHOR1}}\and
  Second Author\affOne\affTwo\orcidlink{\FILL{ORCID\_AUTHOR2}}\lastand\\
  Third Author\affOne\orcidlink{\FILL{ORCID\_AUTHOR3}}%
}
\TIBaffiliations{%
  \affOne \FILL{AFF1}\medskip\\
  \affTwo \FILL{AFF2}\medskip\\
  *Correspondence: \FILL{CORR\_EMAIL}
}

% ---------- Title ----------
\TIBtitle{VLMinators at Sci-ImageMiner 2026 Tasks 1--4: QLoRA Fine-tuning of
  Qwen2.5-VL with Selective Label Disambiguation and Cross-Task Context Injection}

% ---------- Conference ----------
\TIBbundlename[\FILL{BUNDLE\_NAME}]{%
  Sci-ImageMiner 2026: Scientific Image Mining Challenge at ICDAR 2026}
\TIBconferencesession{Sci-ImageMiner 2026 Task Participant Papers}

% ---------- Abstract ----------
\TIBabstract{%
We describe the VLMinators system for the ICDAR 2026 Sci-ImageMiner competition,
in which we participated in all four tasks: figure panel classification, data
extraction, summarization, and visual question answering (VQA) on scientific
figures from atomic layer deposition and etching (ALD/ALE) literature.
For classification (Task~1), we fine-tuned \texttt{Qwen/Qwen2.5-VL-7B-Instruct}
and \texttt{meta-llama/Llama-3.2-11B-Vision-Instruct} using QLoRA (r=32, $\alpha$=64)
and augmented training data with the ACL-Fig and DocFigure external datasets.
The key contribution is \emph{selective label hints}: targeted one-line disambiguation
descriptions injected only for the most confusable class pairs in the 49-class taxonomy,
which improved separation of visually similar chart types such as line vs.\ multiple
line chart and single vs.\ stacked spectra.
A lightweight post-processing pipeline corrected label normalization errors and used
visual inspection to reclassify unknown panels, applying 18 corrections in total.
Our best classification result achieved F1\,=\,0.8059 on the CodaBench evaluation phase
using a six-model weighted ensemble.
For Tasks~2--4, we fine-tuned Qwen2.5-VL-7B-Instruct independently for data extraction
(12,667 samples), summarization (1,573 samples), and VQA (3,041 QA pairs).
A cross-task context-injection pipeline chained outputs sequentially: figure type labels
from Task~1 enriched Task~2 prompts; type and extracted table enriched Task~3;
and all three enriched Task~4 answers, improving numerical specificity without
additional fine-tuning.%
}

\TIBkeywords{scientific figure understanding, vision-language models, QLoRA,
  chart classification, data extraction, summarization, visual question answering,
  ALD/ALE, selective prompting}

\addbibresource{local.bib}

% ---------- Packages ----------
\usepackage{booktabs}
\usepackage{tabularx}
\usepackage{multirow}
\usepackage{graphicx}
\usepackage{xcolor}
\usepackage{listings}
\usepackage{hyperref}

\begin{document}
\maketitle

% ============================================================
\section{Introduction}
% ============================================================

Scientific literature in materials science generates a large volume of complex figures
that encode experimental results in diverse visual formats --- from line charts and
spectra to apparatus diagrams and reaction schemes.
Automatically extracting, summarising, and answering questions about these figures
would significantly accelerate literature analysis and knowledge discovery.

The ICDAR 2026 Sci-ImageMiner competition \cite{sciImageMiner2026} addresses this
challenge through four tasks focused on scientific figures from papers on atomic layer
deposition (ALD) and atomic layer etching (ALE):
\textbf{Task~1} (figure classification into 49 chart/diagram types),
\textbf{Task~2} (extraction of quantitative data as structured tables),
\textbf{Task~3} (generation of concise scientific summaries), and
\textbf{Task~4} (visual question answering).

We present the VLMinators system, which achieved a top-5 finish across all four tasks.
Our approach centers on parameter-efficient fine-tuning of the
Qwen2.5-VL-7B-Instruct vision-language model \cite{qwen25vl} via QLoRA \cite{qlora},
augmented by targeted prompt engineering, lightweight post-processing,
and a cross-task context-injection pipeline.
The main contributions are:
\begin{itemize}
  \item \textbf{Selective label hints} for fine-grained chart disambiguation
        (Section~\ref{sec:task1});
  \item A \textbf{post-processing pipeline} combining rule-based label normalization
        with AI-assisted visual reclassification of unknown panels (Section~\ref{sec:postprocess});
  \item A \textbf{cross-task context-injection framework} that chains task outputs
        to improve downstream tasks without additional training (Section~\ref{sec:context}).
\end{itemize}


% ============================================================
\section{Related Work}
% ============================================================

\paragraph{Scientific figure understanding.}
Chart classification and data extraction from scientific figures have been studied
in the context of document understanding~\cite{docfigure, aclfig}.
Prior work relies on purpose-built CNNs or small transformer models trained on
general chart corpora.
The Sci-ImageMiner competition extends this to a highly specialized domain
(ALD/ALE) with a fine-grained 49-class taxonomy.

\paragraph{Vision-language models.}
Recent large VLMs such as Qwen2.5-VL~\cite{qwen25vl}, LLaMA~\cite{llama32},
LLaVA~\cite{llava15,llava16}, Phi-3.5-Vision~\cite{phi35vision}, and
InstructBLIP~\cite{instructblip} have demonstrated strong visual reasoning.
Fine-tuning via LoRA~\cite{lora} and QLoRA~\cite{qlora} makes adapting 7B+
parameter models tractable on commodity GPUs.

\paragraph{Parameter-efficient fine-tuning.}
QLoRA~\cite{qlora} enables fine-tuning of quantized models by adding low-rank
adaptors to attention layers.
We use the PEFT library~\cite{peft} with r=32, $\alpha$=64 across all our QLoRA
experiments.


% ============================================================
\section{Dataset and Task Overview}
% ============================================================

The competition provides figure images from ALD/ALE papers with panel-level
annotations in JSON format.
Each figure is cropped into individual panels using provided bounding boxes.
The taxonomy for Task~1 contains 49 classes covering data plots (line charts,
scatter plots, spectra, bar charts), scientific diagrams (apparatus, reaction
schemes, process timing), and miscellaneous types (image panels, tables, formulas).

\begin{table}[h!]
  \caption{Dataset split statistics.}\label{tab:data}
  \begin{tabularx}{\linewidth}{lrr}
    \toprule
    \textbf{Split} & \textbf{Figures} & \textbf{Panels} \\
    \midrule
    Train          & ---  & 2,421 \\
    Dev            & ---  & 363   \\
    Test           & 580  & 1,121 \\
    \bottomrule
  \end{tabularx}
\end{table}

For Tasks~2--4 we used additional training files provided by the competition:
12,667 extraction samples, 1,573 summarization samples, and 3,041 VQA pairs.
External data used for Task~1 augmentation: ACL-Fig \cite{aclfig} (3,961~samples)
and DocFigure \cite{docfigure}.


% ============================================================
\section{Approach}
% ============================================================

\subsection{Task 1: Figure Classification}
\label{sec:task1}

\subsubsection{Baselines}

We first evaluated several VLMs in a zero-shot setting to establish baselines:
LLaVA-v1.6-Mistral-7B \cite{llava16}, LLaVA-1.5-7B \cite{llava15},
Phi-3.5-Vision-Instruct \cite{phi35vision}, and InstructBLIP-Vicuna-7B
\cite{instructblip}.
We also fine-tuned three CNN classifiers using the \texttt{timm} library
\cite{timm}: EfficientNet-B0 \cite{efficientnet}, Inception-ResNet-v2
\cite{inceptionresnet}, and SwinV2-Base ($256\times256$) \cite{swinv2}.
CNN models were trained for 30~epochs on competition data augmented with
ACL-Fig, using standard image augmentation and a cross-entropy objective.

\subsubsection{QLoRA Fine-tuning of Qwen2.5-VL and Llama-3.2-Vision}

We fine-tuned \texttt{Qwen/Qwen2.5-VL-7B-Instruct} \cite{qwen25vl} and
\texttt{meta-llama/Llama-3.2-11B-Vision-Instruct} \cite{llama32} using QLoRA
\cite{qlora} (r=32, $\alpha$=64, dropout=0.05, bfloat16).
LoRA adaptors were applied to the attention projection layers
(\texttt{q\_proj}, \texttt{k\_proj}, \texttt{v\_proj}, \texttt{o\_proj}).
We ran multiple experiments (Exp~C--H) varying epochs (3--5), learning rate
($1\times10^{-4}$ vs.\ $2\times10^{-4}$), LoRA rank, and training data
(competition data alone vs.\ augmented with ACL-Fig).
All panels were cropped from figure images using provided bounding boxes and
resized to $448\times448$.

The classification prompt listed all 49 taxonomy labels; the model was
instructed to output only the category name.

\subsubsection{Selective Label Hints}
\label{sec:hints}

Early experiments showed that the model frequently confused visually similar
class pairs, in particular:
single vs.\ multiple line charts,
scatter plots vs.\ multiple scatter plots (especially when a fit curve is present),
single vs.\ stacked vs.\ overlaid spectra, and
process timing diagrams vs.\ other step-function plots.

Rather than adding descriptions for all 49~classes (which degraded performance
on already-correct classes), we injected short disambiguation hints \emph{only}
for the ten most confusable pairs.
For example:
\begin{itemize}
  \item \textit{line chart}: ``SINGLE data series shown as a continuous line/curve''
  \item \textit{scatter plot}: ``x-y data as DISCRETE POINTS or MARKERS; may include a
        single fit/trend curve; SINGLE group or dataset''
  \item \textit{stacked spectra chart}: ``Multiple spectra arranged VERTICALLY STACKED
        (offset along y-axis)''
\end{itemize}
All other classes retained their plain names.
This selective strategy improved dev macro F1 by 3--5 percentage points over
full-description prompts without degrading head-class accuracy.

\subsubsection{Post-processing Pipeline}
\label{sec:postprocess}

Model outputs sometimes differed from the official taxonomy in minor ways
(e.g.\ ``device structure'' instead of ``device structure diagram'').
\texttt{flag\_panels.py} scanned predictions and identified two issue types:
(i)~known label-name mismatches (auto-fixable via a small lookup table) and
(ii)~panels labeled ``unknown.''
\texttt{merge\_classifications.py} then merged corrections back into the
prediction JSON.
For the 15~``unknown'' panels, we used Claude Code to visually inspect each
panel image and assign the correct class --- 11~were reclassified as
\textit{reaction energy profile diagram}, 1~as \textit{reaction scheme}, and
3~as \textit{comparison table}.
In total, 18~corrections were applied.

\subsubsection{Ensemble}
\label{sec:ensemble}

We built weighted majority-vote ensembles using \texttt{src/ensemble.py}.
Weights were set to each model's dev macro F1 score.
The best configuration combined six models: Qwen2.5-VL (Exp~G, 5~epochs),
LLaMA-3.2-11B-Vision (3~epochs), SwinV2-ACLFig, Inception-ResNet-v2,
EfficientNet-B0, and a zero-shot Qwen2.5-VL baseline.
Ensembling with architecturally dissimilar models (VLMs + CNNs) reduced
correlated classification errors on head classes.
Note: we found that ensembling Qwen2.5-VL with LLaMA alone \emph{hurt} performance
due to a high disagreement rate ($\approx$47\%) on fine-grained classes, indicating
that model diversity at the architecture level matters more than diversity from
different fine-tuning runs of the same base model.

% --------------------------------------------------------
\subsection{Task 2: Data Extraction}

We fine-tuned Qwen2.5-VL-7B-Instruct with QLoRA on the 12,667 extraction
training samples.
The model was prompted to produce a Markdown table using axis labels and legend
entries as column headers (with units), outputting only the table without
preamble.
Panels not containing quantitative plots (e.g.\ apparatus diagrams, image panels)
were prompted to output an empty table.

At inference time, the predicted figure type from Task~1 was prepended to the
extraction prompt as \texttt{[Figure type: \textit{label}]}, allowing the model
to set type-appropriate column structure.
For instance, knowing a panel is a ``spectra chart'' guides the model toward
Binding Energy/Intensity columns rather than generic axis labels.

% --------------------------------------------------------
\subsection{Task 3: Summarization}

We fine-tuned Qwen2.5-VL-7B-Instruct with QLoRA on 1,573~summarization
training samples.
The model was instructed to generate 1--3~sentence scientific summaries
describing chart type, axis variables, key trends, and specific numerical
quantities with units.

For inference, we injected both the figure type (Task~1) and the extracted
Markdown table (Task~2) into the summarization prompt:
\begin{verbatim}
[Figure type: line chart]
[Extracted data:
| Temperature (degC) | GPC (A/cycle) |
|---|---|
| 200 | 0.85 |
| 250 | 1.12 |]
Write a concise 1-3 sentence scientific summary...
\end{verbatim}
This grounded the summary in specific numerical values rather than relying
solely on visual interpretation, enabling the model to cite measured quantities.

% --------------------------------------------------------
\subsection{Task 4: Visual Question Answering}

We fine-tuned Qwen2.5-VL-7B-Instruct with QLoRA on 3,041~QA pairs covering
four question types: Yes/No (340), Factoid (1,142), List (201), and
Paragraph (1,358).
The answer type was included in the user prompt so the model could calibrate
response length and format.
We used variable \texttt{max\_new\_tokens} per type: 80 for Yes/No and Factoid,
120 for List, and 300 for Paragraph.

At inference, the full context chain was prepended to each question:
figure type, extracted table, and summary.
This was particularly valuable for Factoid and Paragraph questions that require
citing specific numerical values observable in the extracted table.

% --------------------------------------------------------
\subsection{Cross-Task Context Injection}
\label{sec:context}

\begin{figure}[h!]
  % Replace with an actual TikZ or PDF figure showing the 4-phase chain
  \fbox{\parbox{0.9\linewidth}{
    \centering
    \textbf{[Figure: Context injection pipeline]}\\[4pt]
    Task~1 (classify) $\longrightarrow$ label\\
    $\downarrow$\\
    Task~2 (extract) with \texttt{[Figure type: label]}\\
    $\downarrow$\\
    Task~3 (summarize) with type + table\\
    $\downarrow$\\
    Task~4 (VQA) with type + table + summary
  }}
  \caption{Cross-task context injection pipeline. Each task's output is
    passed as structured text context to the next task's prompt.}
  \label{fig:pipeline}
\end{figure}

The context-injection pipeline (Figure~\ref{fig:pipeline}) requires no
additional fine-tuning and does not increase model parameters.
The key insight is that each downstream task benefits from structured,
text-encoded outputs of earlier tasks: the extraction model benefits from
knowing the chart type; the summarization model benefits from having the
extracted table to cite; and the VQA model benefits from all three.
This is implemented in \texttt{run\_hybrid\_pipeline.py}, which processes all
580~test figures sequentially and generates per-task submission ZIPs.


% ============================================================
\section{Experiments and Results}
% ============================================================

\subsection{Task 1: Classification Results}

Table~\ref{tab:classification} summarises classification performance.
Dev macro F1 measures per-class average (sensitive to tail classes);
dev weighted F1 reflects class-frequency-weighted average.
The CodaBench eval score uses the official competition F1 metric.

\begin{table}[h!]
  \caption{Classification model comparison on dev set and CodaBench eval phase.}\label{tab:classification}
  \begin{tabularx}{\linewidth}{lXrr}
    \toprule
    \textbf{Model} & \textbf{Data} & \textbf{Dev Macro F1} & \textbf{Dev Wtd F1} \\
    \midrule
    Zero-shot LLaVA-v1.6       & —           & —      & —      \\
    Zero-shot Phi-3.5-Vision    & —           & —      & —      \\
    EfficientNet-B0             & Comp.+ACL   & —      & —      \\
    SwinV2-Base                 & Comp.+ACL   & 0.4260 & —      \\
    Inception-ResNet-v2         & Comp.+ACL   & —      & —      \\
    Llama-3.2-11B-Vision        & Competition & 0.5805 & 0.737  \\
    Qwen2.5-VL Exp~G (5~ep)     & Competition & 0.5474 & 0.739  \\
    \midrule
    \textbf{6-model ensemble}   & —           & \textbf{0.568}  & \textbf{0.723}  \\
    \midrule
    \multicolumn{2}{l}{\textit{CodaBench eval phase (F1)}} & \multicolumn{2}{c}{\textbf{0.8059}} \\
    \bottomrule
  \end{tabularx}
\end{table}

The gap between dev macro F1 ($\approx$0.57) and CodaBench F1 (0.8059) reflects
differences in the test set class distribution; macro F1 on dev is strongly
penalized by tail classes with very few examples.

\subsection{Tasks 2--4 Results}

\begin{table}[h!]
  \caption{Results for Tasks 2--4 on the CodaBench test phase.}\label{tab:tasks234}
  \begin{tabularx}{\linewidth}{llX}
    \toprule
    \textbf{Task} & \textbf{Metric} & \textbf{Score} \\
    \midrule
    Task~2 — Extraction    & Weighted / TEDS / RMS         & 40.8 / 64.3 / 17.3 \\
    Task~3 — Summarization & \FILL{TASK3\_METRIC}           & \FILL{TASK3\_SCORE} \\
    Task~4 — VQA           & \FILL{TASK4\_METRIC}           & \FILL{TASK4\_SCORE} \\
    \bottomrule
  \end{tabularx}
\end{table}


% ============================================================
\section{Discussion and Error Analysis}
% ============================================================

\paragraph{Classification.}
Selective label hints were the single most impactful technique for Task~1.
The original scatter plot hint (``NO connecting lines'') caused systematic
errors on panels with a fit curve overlaid on scatter data; replacing it with
a marker-based description (``DISCRETE POINTS or MARKERS; may include a single
fit/trend curve'') fixed the most common confusion.
Ensembling architecturally diverse models (VLMs + CNNs) helped reduce errors on
head classes, but combining VLMs from different families (Qwen vs.\ Llama)
degraded performance due to high pairwise disagreement on fine-grained subtypes.

\paragraph{Extraction.}
The TEDS score of 64.3 indicates that the model recovers table structure
reasonably well but struggles with multi-column layouts, non-standard axis
formats, and panels where data is encoded as colour rather than position.
Context injection (figure type) marginally improved column naming.

\paragraph{Summarization and VQA.}
Injecting extracted tables into summarization prompts improved numerical
specificity, though the model sometimes hallucinated values not present in the
table when the table was empty (non-data panels).
For VQA, the context chain was most beneficial for Paragraph-type questions;
Yes/No and Factoid answers were mostly unaffected by context injection.


% ============================================================
\section{Conclusion}
% ============================================================

We presented the VLMinators system for the ICDAR 2026 Sci-ImageMiner competition.
QLoRA fine-tuning of Qwen2.5-VL-7B-Instruct, combined with selective label
hints and a lightweight post-processing pipeline, achieved F1\,=\,0.8059 on
the classification task eval phase (top-5 ranking).
A cross-task context-injection pipeline chained task outputs sequentially to
improve extraction, summarization, and VQA without additional training.
Future work should explore joint multi-task fine-tuning and better handling of
non-data panels in the extraction and summarization pipelines.


% ============================================================
\section*{Data Availability Statement}
% ============================================================

The competition dataset is provided by the Sci-ImageMiner 2026 organizers and
is available via the CodaBench competition portal.
The ACL-Fig and DocFigure external datasets are publicly available at their
respective repositories \cite{aclfig,docfigure}.
Our code is available at \url{https://github.com/tirtha-v/sci-image-miner}.

% ============================================================
\section*{Underlying and Related Material}
% ============================================================

Code supporting this work is deposited at
\url{https://github.com/tirtha-v/sci-image-miner} (DOI: \FILL{REPO\_DOI}).
The repository includes all fine-tuning scripts, inference pipelines,
ensemble utilities, post-processing tools, and this paper source.
Fine-tuned model adapters will be released on Hugging Face Hub.

% ============================================================
\section*{Author Contributions}
% ============================================================

\textbf{Author~1}: Conceptualization, Methodology, Software (Tasks~2--4,
  hybrid pipeline), Writing --- original draft.\\
\textbf{Author~2}: Methodology, Software (Classification task, post-processing),
  Writing --- review and editing.\\
\textbf{Author~3}: Software (CNN baselines, ensemble), Data curation,
  Writing --- review and editing.

% ============================================================
\section*{Competing Interests}
% ============================================================

The authors declare that they have no competing interests.

% ============================================================
\section*{Funding}
% ============================================================

This work received no specific grant from any funding agency in the public,
commercial, or not-for-profit sectors.

% ============================================================
\section*{Acknowledgements}
% ============================================================

The authors thank the Sci-ImageMiner 2026 organizers for providing the
competition dataset and evaluation platform.

% ============================================================
\section*{References}
% ============================================================

\printbibliography[heading=references]

\end{document}
